\section{Method}
\label{sec:method}

\begin{figure*}[t]
\centering
\resizebox{\textwidth}{!}{%
\begin{tikzpicture}[
  font=\small,
  box/.style={draw, rounded corners=2pt, minimum height=0.9cm, align=center, inner sep=4pt},
  fz/.style={box, fill=frozen!15, draw=frozen, thick},
  tr/.style={box, fill=trained!15, draw=trained, thick},
  tok/.style={draw, minimum width=0.38cm, minimum height=0.42cm, inner sep=0pt},
  arr/.style={-{Stealth[length=2.2mm]}, thick}]
  % inputs and the two streams
  \node[box, fill=gray!10, minimum width=1.9cm] (x) at (0,0) {Meeting\\audio $\bm{x}$};
  \node[fz, minimum width=2.7cm] (enc) at (3.2,1.0) {Audio encoder\\(\S\ref{sec:method-task})};
  \node[fz, minimum width=2.7cm] (dia) at (3.2,-1.0) {Streaming diarizer\\$S{=}8$ slots};
  \node[tr, minimum width=2.3cm] (proj) at (6.5,-1.0) {Speaker projector\\(App.~\ref{app:model})};
  \draw[arr] (x.east) -- ++(0.3,0) |- (enc.west);
  \draw[arr] (x.east) -- ++(0.3,0) |- (dia.west);
  \draw[arr] (dia) -- node[above]{$\bm{P}$} (proj);
  % interleaved sequence
  \foreach \i/\c in {0/frozen,1/frozen,2/frozen,3/frozen,4/frozen,5/spkc,6/frozen,7/frozen,8/frozen,9/frozen,10/frozen,11/spkc} {
    \node[tok, fill=\c!35] (t\i) at (8.9+\i*0.4,0) {};
  }
  \node (dots) at (8.9+12*0.4+0.05,0) {$\cdots$};
  \draw[decorate, decoration={brace, amplitude=4pt}] (t0.north west) -- node[above=4pt, font=\footnotesize]{$\bm{A}_1$: $K{=}5$ acoustic tokens} (t4.north east);
  \node[below=2pt of t5, font=\footnotesize] {$\bm{s}_1$};
  \node[below=2pt of t11, font=\footnotesize] {$\bm{s}_2$};
  \node[font=\footnotesize] at (11.2,-1.25) {interleaved input $\bm{E}$ (Eq.~\ref{eq:interleave})};
  \draw[arr] (enc.east) -| node[pos=0.3, above, font=\footnotesize]{acoustic tokens $\bm{a}_i$} (t8.north);
  \draw[arr] (proj.east) -- node[above, font=\footnotesize]{$\bm{s}_j$} ++(0.75,0) |- (t0.west);
  % LLM
  \node[tr, minimum height=1.9cm, minimum width=2.0cm] (llm) at (15.6,0) {Thinker LLM\\+ LoRA\\(Eq.~\ref{eq:loss})};
  \draw[arr] (dots.east) -- (llm.west);
  % output segments
  \node[box, fill=white, anchor=west] (o1) at (17.25,0.6)
    {$\spk{2}\,\tim{5}$\ \zh{我想知道会议的时间}\ $\tim{48}$};
  \node[box, fill=white, anchor=west] (o2) at (17.25,-0.6)
    {$\spk{3}\,\tim{42}$\ \zh{对，没错。}\ $\tim{56}$};
  \draw[arr] (llm.east) -- ++(0.3,0) |- (o1.west);
  \draw[arr] (llm.east) -- ++(0.3,0) |- (o2.west);
  \node[font=\footnotesize, align=center] at (19.6,1.55) {who, when, what (Eq.~\ref{eq:serial})};
  % talker
  \node[tr, minimum width=1.9cm] (talk) at (24.4,0) {Talker\\voice $\bm{v}_k$};
  \draw[arr] (o1.east) -- ++(0.25,0) |- (talk.west);
  \draw[arr] (o2.east -| o1.east) ++(0,0) -- ++(0.25,0) |- (talk.west);
  \draw[thick] (o2.east) -- (o2.east -| o1.east);
  \node[box, fill=gray!10, minimum width=1.7cm] (y) at (26.9,0) {Translated\\speech};
  \draw[arr] (talk) -- (y);
  % voice prompts
  \draw[arr, dashed, spkc!80!black] (dia.south) -- ++(0,-0.8) -| node[pos=0.2, below, font=\footnotesize]{voice prompt $\bm{v}_k$: solo regions of slot $k$ (Eq.~\ref{eq:alone})} (talk.south);
  % simultaneous loop
  \draw[arr, dotted, trained!80!black] (llm.north) -- ++(0,0.5) -- node[above, font=\footnotesize]{simultaneous: read every $\Delta$\,s, write or $\wait$ (Eqs.~\ref{eq:write}, \ref{eq:bias}); when to write: latent (\S\ref{sec:method-latent})} ++(-6.4,0);
\end{tikzpicture}}
\caption{\sys{}. A frozen audio encoder turns the meeting into acoustic tokens (what is said); a frozen streaming
diarizer gives slot activities $\bm{P}$ that a trained projector turns into speaker tokens $\bm{s}_j$ (who says it).
One speaker token follows every block of $K{=}5$ acoustic tokens. The Thinker writes speaker-attributed, time-stamped
translation segments, offline or read by read, and a Talker speaks each segment in the voice of its speaker slot.
Blue: frozen; orange: trained. The Thinker's gradients reach only the projector, the new token rows and its LoRA
adapters; the Talker is fine-tuned separately on the playback timeline (\S\ref{sec:method-plem}) with the Thinker
fixed. The diarizer and the encoder are never trained.}
\label{fig:arch}
\end{figure*}


\subsection{Task and input: who said what}
\label{sec:method-task}
\label{sec:method-audio}
\label{sec:method-spk}
\label{sec:method-fuse}

\paragraph{Task.} From a single-channel recording $\bm{x}$ of $T$ seconds in which an unknown number of speakers talk,
sometimes at once, the system outputs segments
\begin{equation}
\mathcal{Y}=\big((k_n,\,a_n,\,b_n,\,\bm{y}^{(n)})\big)_{n=1}^{N_\mathcal{Y}},
\label{eq:task}
\end{equation}
where $k_n\in\{1,\dots,S\}$ is the speaker label (arbitrary, but one person keeps one label for the whole meeting),
$a_n<b_n$ the start and end of the source speech and $\bm{y}^{(n)}$ its translation; in speech output $\bm{y}^{(n)}$ is
spoken in the voice of $k_n$. In simultaneous mode segments are produced while $\bm{x}$ is recorded.

\paragraph{Two streams.} The frozen audio encoder of Qwen3-Omni, whose LLM (the ``Thinker'') already translates
single-speaker speech~\citep{xu2025qwen3omni}, maps the recording to acoustic tokens
$\bm{A}=\mathrm{Enc}(\bm{x})=[\bm{a}_1;\dots;\bm{a}_N]\in\R^{N\times d}$ (one per 80\,ms, centred at $c_i$; $d=2048$), which carry no memory of who spoke earlier. A frozen Streaming Sortformer~\citep{medennikov2025streaming}, run in its
streaming mode, keeps every voice in the slot it got when it first appeared and gives slot activities
$\bm{P}=\mathrm{Diar}(\bm{x})\in[0,1]^{F\times S}$  for $F$ frames
of $\delta=80$\,ms and $S=8$ slots; overlapped speech appears as several large entries in one row. We resample $\bm{P}$
to the token centres, $\tilde{\bm{P}}=\mathrm{Interp}(\bm{P};c_{1:N})$ , and summarise every block of $K=5$ rows by a strided
convolution and a two-layer perceptron into one speaker token $\bm{s}_j\in\R^{d}$ (6.3M parameters; the output layer starts at zero, so training
starts from the pretrained Thinker; Appendix~\ref{app:model}). A streaming frame is emitted once and never revised, so a
reader at time $t$ sees what a run on $\bm{x}_{\le t}$ gives (Appendix~\ref{app:causal}).

\paragraph{Interleaved fusion.} Following G-STAR~\citep{peng2026gstar}, each speaker token follows the audio it describes,
\begin{equation}
\bm{E}=[\bm{A}_1;\bm{s}_1;\,\bm{A}_2;\bm{s}_2;\,\dots;\,\bm{A}_M;\bm{s}_M],
\label{eq:interleave}
\end{equation}
with $\bm{A}_j$ the $j$-th block of $K$ acoustic tokens; $\bm{E}$ fills the audio placeholder of the chat template after
an instruction $\bm{c}$ (Appendix~\ref{app:prompt}).

\subsection{Speaker-attributed, time-stamped output}
\label{sec:method-out}
\label{sec:method-train}

\paragraph{Serialisation.} We add 251 time tokens $\tim{0},\dots,\tim{250}$ (80\,ms steps, up to 20\,s) and 8 speaker
tokens $\spk{1},\dots,\spk{8}$ with trainable input and output rows, and write a window's segments speaker first,
ordered by start as in serialized output training~\citep{kanda2020sot},
\begin{equation}
\bm{y}=\mathop{\Vert}_{n=1}^{N_\mathcal{Y}}\big[\spk{k_n}\,\tim{\alpha_n}\,\bm{y}^{(n)}\,\tim{\beta_n}\big],
\label{eq:serial}
\end{equation}
with $\alpha_n,\beta_n$ start and end rounded to 80\,ms. Targets use the diarizer's slots: per meeting, each reference
speaker is mapped to the slot that overlaps most with that speaker's speech over the whole meeting (Hungarian
matching;~\citealp{kuhn1955hungarian}), so the model learns to copy the slot rather than invent a label per window
(Appendix~\ref{app:model}).

\paragraph{Training.} The Thinker $p_\theta(\bm{y}\mid\bm{E},\bm{c})=\prod_u p_\theta(y_u\mid\bm{y}_{<u},\bm{E},\bm{c})$ is trained with the weighted cross-entropy
\begin{equation}
\mathcal{L}=-\frac{\sum_{u}\omega(y_u)\log p_\theta(y_u\mid\bm{y}_{<u},\bm{E},\bm{c})}{\sum_{u}\omega(y_u)},
\label{eq:loss}
\end{equation}
$\omega=2$ for speaker and $1.5$ for time tokens, since one wrong speaker or time token can misattribute a sentence.
Trained are the projector, the new rows and LoRA adapters~\citep{hu2022lora} on all 48 layers (Table~\ref{tab:hyper});
\sys{} (offline) decodes a window at once.

\subsection{Simultaneous interpretation}
\label{sec:method-simul}

\paragraph{Reading and writing.} The model reads every $\Delta=1$\,s, at $t_r=\min(r\Delta,T)$, $r=1,\dots,R$. It sees
$\bm{x}_{\le t_r}$ and a causal speaker stream: the low-latency diarizer finalises a frame only $\lambda=1.04$\,s later,
so later frames are zeroed,
\begin{equation}
P^{(r)}_{f,k}=P_{f,k}\;\mathds{1}\big[f\delta\le t_r-\lambda\big],
\label{eq:causal}
\end{equation}
and $\bm{E}^{(r)}$ is built from $\bm{x}_{\le t_r}$ and $\bm{P}^{(r)}$. With $\hat{\bm{y}}^{(<r)}$ everything committed
before read $r$, the model continues greedily,
\begin{equation}
\hat{\bm{y}}^{(r)}=\mathrm{Greedy}_\theta\big(\cdot\mid\hat{\bm{y}}^{(<r)},\bm{E}^{(r)},\bm{c}\big)\ \text{until}\ \eta_r,
\label{eq:write}
\end{equation}
where the stop token is $\eta_r=\wait$ for $r<R$ and $\eta_R=\eos$; what precedes it is committed and never revised.

\paragraph{Deadline targets.} Each reference translation is split at punctuation into chunks, and each chunk is timed by
the source it covers (word alignments, or the same fraction of the utterance where none exist; Appendix~\ref{app:model}).
A chunk $j$ ending at $e_j$ becomes due one interpreting lag $\epsilon=1$\,s later, at the read
\begin{equation}
r(j)=\min\big(\{r: t_r\ge e_j+\epsilon\}\cup\{R\}\big),
\label{eq:due}
\end{equation}
and the target of read $r$ is the chunks due at $r$, serialised as in Eq.~\eqref{eq:serial}, then $\eta_r$; each read is
one training row whose prefix is the targets of all earlier reads. A deadline model trains new LoRA adapters, the
projector and the new rows on these rows with Eq.~\eqref{eq:loss}, starting from \sys{} (offline). The v9 system,
\sys{}-Simul, has one such model per direction; in v10 one model is trained on both directions (\S\ref{sec:setup}). It is
our strict baseline and the initialisation of latent write alignment.

\subsection{Latent write alignment}
\label{sec:method-latent}

\paragraph{Lattice.} Deadlines fix when each chunk is written; we let the model's likelihood decide within one read. Let
$\mathcal{C}$ be the chunks of a window. A \emph{write plan} $z=(A_{1:R},o_{1:R})$ gives the set
$A_r\subseteq\mathcal{C}$ of chunks written by read $r$, with
\begin{equation}
\{j:r(j)\le r-1\}\subseteq A_r\subseteq\{j:r(j)\le r\},
\label{eq:lattice}
\end{equation}
$A_R=\mathcal{C}$, and each speaker's chunks written in order, so a chunk is written on time or one read late; $o_r$ chooses whether chunks
deferred from read $r-1$ precede the newly due chunks of other speakers. With $Y_r^{o}(A',A)$ the tokens of
$A\setminus A'$ serialised in order $o$ and $\kappa(A')$ the canonical serialisation of the chunks already written, the
posterior over plans is a linear-chain CRF whose states are sets of written chunks and whose potentials are the
Thinker's own log-likelihoods,
\begin{equation}
\begin{aligned}
q_\tau(z)&=\tfrac{1}{Z}\exp(\tau\textstyle\sum_{r}\psi_{r,o_r}(A_{r-1},A_r)),\\
\psi_{r,o}(A',A)&=\log p_\theta(Y_r^{o}\eta_r\mid\bm{E}^{(r)},\kappa(A')),
\end{aligned}
\label{eq:post}
\end{equation}
with $A_0=\emptyset$, $Y_r^o=Y_r^o(A',A)$ and the instruction $\bm{c}$ omitted.
Since the history depends only on the set $A'$, the plan score factorises over reads, and $Z$ and every edge marginal
$q_\tau(r,A'\!\to\!A,o)$ are exact by forward--backward over the feasible sets.

\paragraph{EM.} The E-step scores every feasible transition with the current model; the M-step minimises the
posterior-weighted form of Eq.~\eqref{eq:loss} with $q_\tau$ held fixed,
\begin{equation}
\mathcal{L}_{\mathrm{W}}=-\frac{1}{D}\sum_{e}q_\tau(e)\,
\log_\omega p_\theta\big(Y_r^{o}\eta_r\mid\cdot\big),
\label{eq:mstep}
\end{equation}
where $e=(r,A'\!\to\!A,o)$ ranges over the lattice edges, $\log_\omega$ applies the token weights and $D$ the normaliser of Eq.~\eqref{eq:loss}. With all mass on the deadline
plan this is the deadline loss exactly; we anneal with $\tau=0.5$ (Appendix~\ref{app:crf}).

\paragraph{Write bias.} One more read makes the same translation more likely, so EM moves mass to the later write and
the decoder learns to wait longer (\S\ref{sec:res-drift}). At decoding we add one constant to the wait logit at every
step of every read but the last,
\begin{equation}
\tilde{\ell}_u(\wait)=\ell_u(\wait)+\bias,\qquad \bias=-1 .
\label{eq:bias}
\end{equation}
$\bias$ is a fixed decoding constant like greedy search or $\Delta$, chosen once on five development meetings as the
value whose text lag is closest to the deadline model's, for both directions (Appendix~\ref{app:calib}).

\subsection{Speaking in each speaker's voice}
\label{sec:method-voice}

\paragraph{Voices.} Slot $k$'s solo frames are
\begin{equation}
m_{f,k}=\mathds{1}\big[P_{f,k}>0.5\big]\prod_{k'\neq k}\mathds{1}\big[P_{f,k'}<0.5\big];
\label{eq:alone}
\end{equation}
runs of at least 0.4\,s are concatenated into a voice prompt $\bm{v}_k$ of at most 8\,s, in simultaneous mode only from
audio heard before the slot's first commitment. The UniSS Talker~\citep{uniss2025}, which predicts BiCodec
tokens~\citep{wang2025spark}, speaks each slot's segments with prompt $\bm{v}_k$. Each speaker has a stream of its own;
within a stream, segment $n$ of duration $D_n$ starts at
\begin{equation}
\pi_n=\max\big(\zeta_n,\;\pi_{n-1}+D_{n-1}\big),
\label{eq:layout}
\end{equation}
with $\zeta_n$ its release time (the source start offline). The baseline Talker plays every committed block in full at a
fixed rate, so a stream that receives more speech than its source builds a queue.

\subsection{Speaking on the playback timeline}
\label{sec:method-plem}

\paragraph{Lattice.} A stream's blocks $i=1,\dots,I$ have $W_i$ target words and are released at $\nu_i$
(\S\ref{sec:method-online}). For each block the Talker proposes candidates $c\in C_i$ (8 samples at two speaking rates)
of duration $D(c)$. The backlog $h_i$, how long block $i$ waits for earlier speech, is rounded up ($\lceil\cdot\rceil$) to an 80\,ms grid
up to 8\,s:
$h_{i+1}=T_i(h_i,c_i)$ with
\begin{equation}
T_i(h,c)=\big\lceil\max(0,\,\nu_i+h+D(c)-\nu_{i+1})\big\rceil,
\label{eq:backlog}
\end{equation}
and with a block's words spread evenly over its playback they are heard $L_i(h,c)=W_ih+\tfrac{W_i+1}{2}D(c)$
word-seconds after $\nu_i$ in total. A linear-chain CRF over $(h_i,c_i)$ has potentials
\begin{equation}
\phi_i(h,c)=0.5\,\bar{\ell}(c)+5\,f(c)-\lambda L_i(h,c),
\label{eq:plpot}
\end{equation}
with $\bar{\ell}$ the Talker's mean token log-probability and $f$ the negative CTC loss per character under an
independent recogniser; forward--backward over the backlog is exact.

\paragraph{Latency budget and training.} $\lambda$ is the dual variable of $\min_q\mathrm{KL}(q\Vert q_0)$ subject to
$\mathbb{E}_q[\sum_iL_i]/\sum_iW_i\le b$, a posterior-regularised problem~\citep{ganchev2010posterior} whose solution is the
CRF of Eq.~\eqref{eq:plpot}; the expected delay falls monotonically in $\lambda$, so bisection finds $\lambda^\star$. The
budget $b$ is the rule-based Talker's own delay minus 0.3\,s. The Talker is fine-tuned on each block's candidates weighted
by their posterior given the backlog the rule-based path reached, which brings the backlog cost of later blocks into the
current one through the backward messages (\emph{PL-EM}; Appendix~\ref{app:plem}). At inference it samples two
candidates per block at a backlog-dependent rate and keeps the one with the better CTC fidelity.

\subsection{Running on a whole meeting}
\label{sec:method-online}

\paragraph{Online windows.} A window opens at the first frame where any slot of the low-latency diarizer is active and
closes after 0.6\,s with no active slot or at 20\,s (v9: 1\,s). All audio is processed. A window's opening and end are
known only $\lambda$ after the frames that decide them, so commits are delayed to the opening, and those of the last read
to the moment the end is known, with frames timed as the diarizer emits them, in blocks of nine (\emph{chunk protocol},
Appendix~\ref{app:chunk}). Slots persist across windows; the decoder restarts at each window. Blocks of one speaker
committed in the same read are merged before they reach the Talker (\emph{read-level merging}, rl).
